\section{Алгебры с двумя операциями. Кольца, поля (3 теоремы). Области целостности. Примеры.}

\subsection*{Кольца}

\paragraph{Определение}
\textbf{Кольцо} — это множество $M$ с двумя бинарными операциями $+$ (сложение) и $\ast$ (умножение), удовлетворяющее следующим аксиомам:
\begin{enumerate}
    \item \textbf{Абелева группа по сложению:} ассоциативность, коммутативность, существование нуля ($0$) и противоположного элемента ($-a$).
    \item \textbf{Полугруппа по умножению:} ассоциативность ($a \ast (b \ast c) = (a \ast b) \ast c$).
    \item \textbf{Дистрибутивность} умножения относительно сложения: 
    $a \ast (b + c) = a \ast b + a \ast c$ и $(a + b) \ast c = a \ast c + b \ast c$.
\end{enumerate}

\paragraph{Виды колец}
\begin{itemize}
    \item \textbf{Коммутативное:} если $a \ast b = b \ast a$.
    \item \textbf{Кольцо с единицей:} если существует $1 \in M$, такой что $a \ast 1 = 1 \ast a = a$ (моноид по умножению).
\end{itemize}

\paragraph{Теорема 1 (Свойства кольца)}
В любом кольце выполняются соотношения:
1. $0 \ast a = a \ast 0 = 0$.
2. $a \ast (-b) = (-a) \ast b = -(a \ast b)$.
3. $(-a) \ast (-b) = a \ast b$.
4. $-a = a \ast (-1)$ (в кольце с единицей).

\textit{Доказательство (пункт 1):}
$0 \ast a = (0 + 0) \ast a = 0 \ast a + 0 \ast a$. Прибавив к обеим частям $-(0 \ast a)$, получим $0 = 0 \ast a$. $\square$

\subsection*{Области целостности}

\paragraph{Делители нуля}
Если в кольце для ненулевых $x, y$ выполняется $x \ast y = 0$, то $x$ и $y$ называются \textbf{делителями нуля}. 
\textit{Пример:} В машинной арифметике $(\mathbb{Z}_{2^{16}})$ имеем $256 \ast 256 = 2^{16} = 0$, хотя $256 \neq 0$.

\paragraph{Определение}
\textbf{Область целостности} — это коммутативное кольцо с единицей, не имеющее делителей нуля.
\textit{Пример:} $(\mathbb{Z}; +, \ast)$ — область целостности.

\paragraph{Теорема 2 (Закон сокращения)}
В кольце закон сокращения ($a \ast b = a \ast c \implies b = c$ при $a \neq 0$) выполняется тогда и только тогда, когда в кольце нет делителей нуля.

\textit{Доказательство:}
$a \ast b = a \ast c \iff a \ast b - a \ast c = 0 \iff a \ast (b - c) = 0$.
Если делителей нуля нет и $a \neq 0$, то $b - c = 0 \implies b = c$.
Если есть делитель нуля $x \ast y = 0 \ (x, y \neq 0)$, то $x \ast y = x \ast 0$, но $y \neq 0$, т.е. сокращение невозможно. $\square$

\subsection*{Поля}

\paragraph{Определение}
\textbf{Поле} — это коммутативное кольцо с единицей ($1 \neq 0$), в котором каждый ненулевой элемент обратим по умножению. 
\textit{Образно говоря:} Поле — это абелева группа по сложению и $M \setminus \{0\}$ — абелева группа по умножению.

\paragraph{Теорема 3 (Свойство поля)}
Всякое поле является областью целостности (т.е. не имеет делителей нуля).

\textit{Доказательство:}
Пусть $a \ast b = 0$ и $a \neq 0$. Так как мы в поле, существует $a^{-1}$. 
Умножим слева: $a^{-1} \ast (a \ast b) = a^{-1} \ast 0 \implies (a^{-1} \ast a) \ast b = 0 \implies 1 \ast b = 0 \implies b = 0$. 
Значит, произведение ненулевых элементов не может быть нулем. $\square$

\subsection*{Примеры}
\begin{enumerate}
    \item $(\mathbb{Z}; +, \ast)$ — коммутативное кольцо, область целостности, но \textbf{не поле} (нет обратных).
    \item $(\mathbb{Q}; +, \ast)$, $(\mathbb{R}; +, \ast)$ — поля.
    \item $(\mathbb{Z}_n; +, \ast)$ — кольцо вычетов. Является полем тогда и только тогда, когда $n$ — простое число.
    \item Машинная арифметика $(\mathbb{Z}_{2^{16}})$ — кольцо с делителями нуля (не область целостности).
\end{enumerate}